\documentclass[sigconf,nonacm]{acmart}

\usepackage{booktabs}
\usepackage{multirow}
\usepackage{graphicx}

\settopmatter{printacmref=false, printccs=false, printfolios=true}
\setcopyright{none}

\begin{document}

% DRAFT v1 (11/14). Basic parts filled; method/results of our own approach are being developed.
\title{Node Classification from Small Citation Graphs to a Large Product Graph}

\author{Shunyang Huang}
\author{Mingyu Lin}
% TODO: affiliation (same as cv-3Drecon report).

\begin{abstract}
We study semi-supervised node classification on two small citation graphs, Cora and Citeseer, and on the large co-purchase graph ogbn-products.
We build a single code base that covers feature-only, label-only and graph neural network models, and evaluates all of them under the standard splits with the official evaluator.
On top of these models we adapt Cooperation of Experts (CoE), a large-margin ensemble originally proposed for heterogeneous multiplex networks, to node classification, and combine it with ensemble self-training.
On ogbn-products this improves its best expert by 1.6 points and the best reference model with Correct and Smooth by 0.5 points, mainly by learning how to combine label propagation with a feature-based model; on the citation graphs, where all experts are GNNs that make similar errors, it only matches the best expert.
\end{abstract}

\maketitle

\section{Introduction}

Node classification asks for the label of every node in a graph when only a subset of nodes is labelled.
The node features alone are often weak predictors: in a citation graph, the bag-of-words of a paper tells part of its topic, and the papers it cites tell much of the rest.
Graph neural networks (GNNs) combine both sources by propagating features along edges~\cite{kipf2017gcn,velickovic2018gat,gasteiger2019appnp,hamilton2017sage}, while classical label propagation uses the graph and the known labels only~\cite{zhu2002lp,zhou2004lgc}.

The course task contains two settings that stress different parts of a method.
Cora and Citeseer~\cite{sen2008collective} have a few thousand nodes and only 20 labelled nodes per class in the standard split~\cite{yang2016planetoid}, so the main difficulty is learning from very few labels.
ogbn-products~\cite{hu2020ogb} has 2.4 million nodes and 62 million edges, so a method must also scale, and its split by sales rank makes the test nodes differ from the training nodes.

We first build a common pipeline and evaluate seven reference models on all three datasets under the same protocol (Section~\ref{sec:baselines}).
We then adapt Cooperation of Experts (CoE)~\cite{wang2025coe} to node classification.
CoE combines several experts through a learned cooperation matrix trained with a large-margin objective.
On graphs, the experts are models that use different information: node features only, labels only, or both through message passing.
Because the experts fit the few training labels almost perfectly, we fit the cooperation matrix on out-of-fold predictions, and we use the combined prediction to produce pseudo-labels for a round of self-training (Section~\ref{sec:method}).
CoE helps on ogbn-products, where label propagation and a feature-based model make different errors, and matches the best expert on the citation graphs, also when they have many labels (Section~\ref{sec:coe-results}).

\section{Related Work}

\paragraph{Graph neural networks.}
GCN~\cite{kipf2017gcn} averages the transformed features of a node and its neighbours with a symmetric normalisation.
GAT~\cite{velickovic2018gat} learns attention weights over neighbours.
APPNP~\cite{gasteiger2019appnp} separates prediction from propagation and propagates the predictions with personalised PageRank, which allows a large receptive field without many layers.
GraphSAGE~\cite{hamilton2017sage} keeps a separate weight for the node itself and for the mean of its neighbours.
SIGN~\cite{rossi2020sign} precomputes the propagated features $\hat{A}^k X$ for several hops and trains an MLP on them, which makes training on large graphs as cheap as training an MLP.

\paragraph{Label propagation.}
Label propagation~\cite{zhu2002lp,zhou2004lgc} spreads the known labels along edges and relies on the assumption that neighbouring nodes share labels.
Correct and Smooth (C\&S)~\cite{huang2021cs} applies two propagation steps on the output of any base model: it first propagates the errors on the training nodes to correct the predictions, and then smooths the corrected predictions over the graph.
With a simple base model it matches or exceeds many GNNs on OGB datasets.

\paragraph{Self-training on graphs.}
Li et al.~\cite{li2018deeper} showed that GCNs need many labels to propagate information over the whole graph, and proposed to add confident predictions as pseudo-labels, including a variant that only keeps nodes on which the GCN and a random-walk model agree.
Liu et al.~\cite{liu2022drgst} showed that high-confidence pseudo-labels can still shift the training distribution, and reweighted them accordingly.

\paragraph{Ensembles.}
Bagging~\cite{breiman1996bagging} averages models trained on resampled data, and stacking~\cite{wolpert1992stacking} trains a combiner on held-out predictions of the base models.
Cooperation of Experts~\cite{wang2025coe} encodes each semantic space of a heterogeneous multiplex network with a dedicated expert, and combines the experts with a cooperation matrix trained with a large-margin loss; it is designed for multi-relational data.
Here we apply the same combination to single-relation graphs, where the experts differ in the information and propagation scheme they use.

\section{Task and Data}

Given a graph $G=(V,E)$ with $n$ nodes, node features $X\in\mathbb{R}^{n\times d}$ and labels $y_i\in\{1,\dots,C\}$ for the training nodes $V_L\subset V$, the goal is to predict the labels of the remaining nodes.
The setting is transductive: the whole graph, including the features of the test nodes, is available during training, but their labels are not.

Table~\ref{tab:data} lists the three datasets.
For Cora and Citeseer we use the public split of Yang et al.~\cite{yang2016planetoid} through PyTorch Geometric~\cite{fey2019pyg}: 20 labelled nodes per class, 500 validation nodes and 1{,}000 test nodes.
Features are binary bag-of-words vectors, which we normalise to unit row sum.
Citeseer has 48 isolated nodes and several small connected components without any labelled node.
For ogbn-products we use the official split of OGB~\cite{hu2020ogb}: nodes are Amazon products, edges connect products bought together, features are 100-dimensional reduced bag-of-words of the product descriptions, and labels are 47 product categories.
The split is by sales rank: the 8\% most popular products are used for training, the next 2\% for validation and the rest for testing.
All graphs are treated as undirected without self loops.

\begin{table*}[t]
\caption{Datasets and splits.}
\label{tab:data}
\begin{tabular}{lrrrrr}
\toprule
Dataset & Nodes & Edges & Features & Classes & Train / Val / Test \\
\midrule
Cora & 2{,}708 & 5{,}278 & 1{,}433 & 7 & 140 / 500 / 1{,}000 \\
Citeseer & 3{,}327 & 4{,}552 & 3{,}703 & 6 & 120 / 500 / 1{,}000 \\
ogbn-products & 2{,}449{,}029 & 61{,}859{,}140 & 100 & 47 & 196{,}615 / 39{,}323 / 2{,}213{,}091 \\
\bottomrule
\end{tabular}
\end{table*}

We report classification accuracy on the test nodes.
For ogbn-products we use the official OGB evaluator.
All model choices, including early stopping and the hyper-parameters of post-processing, are made on validation accuracy; test labels are used only for the reported numbers.

\section{Methods}
\label{sec:method}

\subsection{Base models}

We use $\hat{A}=\tilde{D}^{-1/2}(A+I)\tilde{D}^{-1/2}$ for the symmetrically normalised adjacency with self loops.
All propagation steps are sparse matrix products, so no model materialises per-edge messages and all of them run full-graph on ogbn-products.

\begin{itemize}
\item \textbf{MLP} uses the node features only.
\item \textbf{Label propagation (LP)} uses the graph and the training labels only. With $Y^{(0)}$ the one-hot matrix of training labels (zero rows elsewhere) it iterates $Y^{(t+1)}=(1-\alpha)Y^{(0)}+\alpha\hat{A}Y^{(t)}$ for 50 steps with $\alpha=0.9$. A node that receives no label mass, because its connected component has no labelled node, has no prediction.
\item \textbf{GCN}~\cite{kipf2017gcn}, \textbf{GAT}~\cite{velickovic2018gat}, \textbf{APPNP}~\cite{gasteiger2019appnp} and \textbf{GraphSAGE}~\cite{hamilton2017sage} with mean aggregation, trained full-batch. GAT is used on the small graphs only.
\item \textbf{SIGN}~\cite{rossi2020sign} with a linear encoder for each of the precomputed features $\hat{A}^k X$, $k=0,\dots,K$, followed by an MLP; it is trained with mini-batches of nodes.
\end{itemize}

\paragraph{Correct and Smooth.}
We apply C\&S~\cite{huang2021cs} to the output of every model, propagating the training labels only.
The correction and smoothing strengths are chosen on validation accuracy from $\{0, 0.5, 0.8, 0.95\}$; on ogbn-products the largest correction strength is 1.0.
Strength 0 switches a step off, so C\&S is not applied when it does not help.

\subsection{Cooperation of Experts for node classification}

% DRAFT: this section describes the current implementation and will be revised with the results.
Let $p_{j,i}\in\mathbb{R}^C$ be the predicted class distribution of expert $j\in\{1,\dots,K\}$ on node $i$.
Following CoE~\cite{wang2025coe}, the experts are combined by a cooperation matrix,
\begin{equation}
z_i=\sum_{j=1}^{K}\Theta_j\,p_{j,i},\qquad \hat{y}_i=\operatorname{softmax}(z_i),
\end{equation}
with $\Theta_j\in\mathbb{R}^{C\times C}$, so that the trust in an expert can depend on the class it predicts.
$\Theta$ is trained with cross-entropy and the large-margin term of CoE, which asks the probability of the true class to exceed a smooth maximum of the other classes:
\begin{equation}
\mathcal{L}(\Theta)=\frac{1}{|V_L|}\sum_{i\in V_L}\Big[-\log \hat{y}_{i,y_i}-\gamma\big(\hat{y}_{i,y_i}-m_i\big)\Big],
\end{equation}
where $m_i=\frac{1}{\beta}\log\sum_{c\neq y_i}\exp(\beta\,\hat{y}_{i,c})$ is a smooth maximum over the other classes, with $\beta=100$ as in CoE.
We initialise $\Theta_j=s\,a_j I$, where $a_j$ is the accuracy of expert $j$, which corresponds to accuracy-weighted voting, and add an $\ell_2$ penalty $\lambda\|\Theta-\Theta^{(0)}\|^2$.
$\gamma$ and $\lambda$ are chosen on validation accuracy.

\paragraph{Out-of-fold fitting.}
In the original setting, $\Theta$ is fitted on the predictions of the experts on their own training nodes.
In semi-supervised node classification this does not work: with 20 labels per class every expert fits its training nodes almost perfectly, and label propagation returns the training labels unchanged.
We therefore fit $\Theta$ on out-of-fold predictions.
The training nodes are split into $F$ class-balanced folds; for each fold every expert is trained on the other folds and predicts the held-out fold.
The fitted $\Theta$ is then applied to experts trained on all training nodes.

\paragraph{Experts.}
The experts are chosen to differ in the information they use: MLP (features), LP (labels and graph), and GNNs with different propagation schemes (GCN, GAT, APPNP, GraphSAGE, SIGN).
On ogbn-products we use MLP, LP and SIGN; the full-batch GNNs are left out because the out-of-fold loop would train each of them $F+1$ times per round.

\paragraph{Ensemble self-training.}
The combined prediction is used to select pseudo-labels: unlabelled nodes outside the validation set whose combined confidence is at least $\tau$ receive their predicted label.
All experts are retrained on the training labels and the pseudo-labels, with the two losses averaged separately, $\mathcal{L}=\mathcal{L}_{\text{true}}+w\,\mathcal{L}_{\text{pseudo}}$, so that a large number of pseudo-labels does not outweigh the true labels.
The out-of-fold folds are still formed over the true training labels only.

\section{Experimental Setup}

\paragraph{Implementation.}
All models are implemented in PyTorch with PyTorch Geometric~\cite{fey2019pyg}; ogbn-products is loaded and evaluated with the OGB package~\cite{hu2020ogb}.
Experiments run on NVIDIA A100 40\,GB GPUs.

\paragraph{Hyper-parameters.}
On Cora and Citeseer all models use 64 hidden units (GAT: 8 heads of 8 units), dropout 0.5 (GAT: 0.6), learning rate 0.01 (GAT: 0.005), weight decay $5\times10^{-4}$, Adam, up to 500 epochs and early stopping with patience 100 on validation accuracy.
APPNP uses 10 propagation steps with teleport probability 0.1, and SIGN uses 3 hops.
On ogbn-products the full-batch models use three layers of 256 units with batch normalisation and learning rate 0.01; SIGN uses 5 hops, 512 hidden units, input dropout 0.2, learning rate 0.001 and batches of 50{,}000 nodes; the MLP uses three layers and learning rate 0.001.
Self-training uses $\tau=0.7$ on the small graphs and $\tau=0.9$ on ogbn-products, and CoE uses $F=5$ folds on the small graphs and $F=3$ on ogbn-products.
All hyper-parameters are listed in \texttt{configs/default.yaml} of our code.

\paragraph{Seeds.}
We report mean and standard deviation over 10 seeds on Cora and Citeseer, where the test set has only 1{,}000 nodes and the standard deviation across seeds is 0.5--1.0 points.
On ogbn-products the test set has 2.2 million nodes and the standard deviation is 0.1--0.3 points, so we use 3 seeds for the reference models (5 for the MLP) and a single seed for the CoE runs, which take about three hours each.
Label propagation is deterministic.

\section{Results}

\subsection{Reference models}
\label{sec:baselines}

\begin{table*}[t]
\caption{Test accuracy (\%) of the reference models, without and with C\&S post-processing. Mean $\pm$ standard deviation over seeds; best per dataset in bold. GAT is not run on ogbn-products.}
\label{tab:baselines}
\begin{tabular}{l cc cc cc}
\toprule
 & \multicolumn{2}{c}{Cora} & \multicolumn{2}{c}{Citeseer} & \multicolumn{2}{c}{ogbn-products} \\
\cmidrule(lr){2-3}\cmidrule(lr){4-5}\cmidrule(lr){6-7}
Model & plain & +C\&S & plain & +C\&S & plain & +C\&S \\
\midrule
MLP & 58.4\,$\pm$\,0.9 & 80.9\,$\pm$\,0.9 & 60.0\,$\pm$\,0.7 & 70.1\,$\pm$\,0.5 & 60.3\,$\pm$\,0.1 & 77.9\,$\pm$\,0.1 \\
Label propagation & 71.0 & -- & 49.8 & -- & 71.3 & -- \\
GCN & 82.0\,$\pm$\,0.5 & 82.1\,$\pm$\,0.7 & 71.2\,$\pm$\,0.5 & 71.4\,$\pm$\,0.8 & -- & -- \\
GAT & 82.3\,$\pm$\,1.0 & 82.4\,$\pm$\,1.1 & \textbf{72.0\,$\pm$\,0.7} & 72.0\,$\pm$\,0.7 & -- & -- \\
APPNP & \textbf{83.2\,$\pm$\,0.9} & \textbf{83.2\,$\pm$\,0.9} & 71.1\,$\pm$\,0.7 & 71.1\,$\pm$\,0.7 & -- & -- \\
GraphSAGE & 81.3\,$\pm$\,0.4 & 82.9\,$\pm$\,0.6 & 69.5\,$\pm$\,0.8 & 70.1\,$\pm$\,1.0 & 78.8\,$\pm$\,0.3 & \textbf{79.0\,$\pm$\,0.4} \\
SIGN & 72.5\,$\pm$\,1.9 & 80.7\,$\pm$\,1.0 & 63.2\,$\pm$\,1.4 & 69.8\,$\pm$\,1.0 & 77.7\,$\pm$\,0.1 & 78.7\,$\pm$\,0.1 \\
\bottomrule
\end{tabular}

\end{table*}

Table~\ref{tab:baselines} shows the reference models.
Our GNN results on the small graphs are close to the published ones: GCN reaches 82.0\% on Cora and 71.2\% on Citeseer (81.5\% and 70.3\% reported in~\cite{kipf2017gcn}), GAT 82.3\% and 72.0\% (83.0\% and 72.5\% in~\cite{velickovic2018gat}), and APPNP is the best model on Cora with 83.2\%.
On ogbn-products, the MLP reaches 60.3\% (61.1\% on the OGB leaderboard), GCN 77.6\% and GraphSAGE 78.8\%.
Our SIGN reaches 78.0\%, below the 80.5\% reported for SIGN on the leaderboard; longer training and smaller batches did not close this gap.

Three observations are relevant for the rest of the report.
First, most of the gap between feature-only and graph models comes from label smoothness alone: C\&S on top of an MLP raises its accuracy from 58.4\% to 80.9\% on Cora and from 60.3\% to 77.9\% on ogbn-products.
On the small graphs, C\&S adds little to the GNNs, which already propagate information over the graph.
Second, label propagation is limited by coverage: on Citeseer it reaches only 49.8\%, because the isolated nodes and the components without labelled nodes receive no label mass.
Third, on ogbn-products validation accuracy is about 93\% while test accuracy is about 78\%, for every model.
The test products have lower sales rank than the training and validation products, so validation accuracy overestimates test accuracy, and model selection on validation is conservative.

\subsection{Cooperation of Experts}
\label{sec:coe-results}

\begin{table}[t]
\caption{Test accuracy (\%) of CoE against its experts and simpler combinations. The expert in parentheses is the one with the best validation accuracy, which is the model one would deploy; the best expert on test is shown for reference. ogbn-products: one seed. $^\dagger$One round of self-training (see text); on the small graphs the number of rounds is chosen on validation.}
\label{tab:coe}
\resizebox{\columnwidth}{!}{\begin{tabular}{l ccc}
\toprule
Method & Cora & Citeseer & ogbn-products \\
\midrule
\multicolumn{4}{l}{\textit{Without self-training}} \\
Best expert, chosen on validation & 82.3\,$\pm$\,1.0 (GAT) & 71.1\,$\pm$\,0.7 (APPNP) & 77.9 (SIGN) \\
Best expert on test (reference) & 83.2\,$\pm$\,0.9 (APPNP) & 72.0\,$\pm$\,0.7 (GAT) & 77.9 (SIGN) \\
Mean voting & 82.8\,$\pm$\,0.4 & 71.3\,$\pm$\,0.6 & 79.1 \\
Accuracy-weighted voting & 82.9\,$\pm$\,0.5 & 71.6\,$\pm$\,0.6 & 78.9 \\
Stacking ($\gamma=0$) & 83.0\,$\pm$\,0.5 & 71.5\,$\pm$\,0.7 & 79.6 \\
CoE & 82.7\,$\pm$\,0.5 & 71.6\,$\pm$\,0.6 & 79.5 \\
\midrule\multicolumn{4}{l}{\textit{With self-training}} \\
Single-model self-training & 82.7\,$\pm$\,0.8 (GAT) & 72.8\,$\pm$\,1.2 (APPNP) & 79.4 (SIGN$^\dagger$) \\
Best expert, ensemble pseudo-labels & 83.7\,$\pm$\,0.4 (GAT) & 73.1\,$\pm$\,0.7 (GAT) & 80.5 (SIGN) \\
CoE, ensemble pseudo-labels & 83.7\,$\pm$\,0.4 & 72.3\,$\pm$\,0.4 & 80.8 \\
\bottomrule
\end{tabular}
}
\end{table}

Table~\ref{tab:coe} compares CoE with its experts and with simpler ways of combining them.

\paragraph{ogbn-products.}
CoE reaches 79.5\%, 1.6 points above its best expert (SIGN, 77.9\%), 0.7 points above SIGN with C\&S (78.8\%) and 0.5 points above the best reference result (GraphSAGE with C\&S, 79.0\%), while the standard deviation across seeds on this dataset is 0.1--0.2 points.
Learned combination is clearly better than voting (79.6\% for stacking and 79.5\% for CoE against 78.9\% for accuracy-weighted voting): the three experts differ strongly in accuracy, and the cooperation matrix learns how much to trust each of them for each class.
Most of the gain comes from the label propagation expert.
Combining only the MLP and SIGN gives 78.3\%, 0.2 points above SIGN.
SIGN uses the features of up to five hops but no labels, and label propagation uses the labels but no features, so the two make different errors; CoE learns to combine them and does so better than C\&S, the usual way of adding label propagation to a feature-based model.

\paragraph{Cora and Citeseer.}
Here CoE performs like voting.
It is 0.4--0.6 points above the expert that validation would choose, but below the expert that happens to be best on test (APPNP on Cora, GAT on Citeseer), and the differences between the combination methods are within one standard deviation.
On Citeseer the validation search picks the strongest pull towards the voting initialisation in every seed, so the cooperation matrix barely moves.

\paragraph{Is the number of labels the reason?}
To separate the effect of the label budget from that of the experts, we repeat the comparison on a random 60/20/20 split of Cora and Citeseer, which gives 1{,}624 and 1{,}996 training labels instead of 140 and 120 (Table~\ref{tab:random60}).
With many labels the validation search no longer constrains the cooperation matrix, but CoE still only matches the best experts.
The number of labels is therefore not the main reason; the experts are.
On the citation graphs, every GNN already propagates information along the same edges and is trained end-to-end on the labels, so the experts make similar errors, and label propagation (83.0\% on Cora with 60\% labels) is weaker than all GNNs without adding new information.
On ogbn-products, the strongest expert does not use the labels at all, which leaves room for label propagation to correct it.

\begin{table}[t]
\caption{CoE on Cora and Citeseer with a random 60/20/20 split (10 seeds, no self-training).}
\label{tab:random60}
\resizebox{\columnwidth}{!}{\begin{tabular}{l cc}
\toprule
Method & Cora & Citeseer \\
\midrule
Best expert, chosen on validation & 86.4\,$\pm$\,0.5 (SAGE) & 76.4\,$\pm$\,0.2 (APPNP) \\
Best expert on test (reference) & 87.0\,$\pm$\,0.2 (APPNP) & 76.9\,$\pm$\,0.8 (SAGE) \\
Accuracy-weighted voting & 86.4\,$\pm$\,0.4 & 77.4\,$\pm$\,0.4 \\
Stacking ($\gamma=0$) & 86.5\,$\pm$\,0.3 & 77.5\,$\pm$\,0.5 \\
CoE & 86.5\,$\pm$\,0.3 & 77.5\,$\pm$\,0.7 \\
\bottomrule
\end{tabular}
}
\end{table}

\paragraph{Large-margin term.}
The large-margin term of CoE does not help in our setting: stacking with $\gamma=0$ is within 0.3 points of CoE on all datasets and splits.
In the original setting of CoE the experts see different semantic spaces of a multiplex network; here they all predict from the same graph, and with out-of-fold fitting the cross-entropy alone is sufficient.

\paragraph{Ensemble self-training.}
Using the combined prediction to select pseudo-labels raises every expert: on Cora APPNP improves from 83.2\% to 84.3\%, on Citeseer GAT from 72.0\% to 73.1\%, and on ogbn-products SIGN from 77.9\% to 80.5\%, with CoE reaching 83.7\%, 72.3\% and 80.8\%.
The selected pseudo-labels are 84\%, 72\% and 91\% correct, for about 2{,}000, 2{,}500 and 1.7 million nodes.
On the small graphs, however, single-model self-training gives the same improvement (e.g.\ APPNP 84.0\% on Cora, GCN 73.5\% on Citeseer), so the benefit there comes from self-training itself and not from the ensemble.
On ogbn-products the ensemble pseudo-labels are better: after one round, SIGN trained on its own pseudo-labels reaches 79.4\%, and on the CoE pseudo-labels 80.5\%; a second round of single-model self-training reaches 80.0\%.
On this dataset validation accuracy does not follow test accuracy across rounds (it stays at about 93\% while test accuracy rises by 1.5--2.5 points), because the validation products come from the same high sales ranks as the training products.
Model selection on validation would therefore always discard self-training on ogbn-products; we report the fixed single round used by ensemble self-training for both methods.

\section{Discussion and Limitations}

Our SIGN on ogbn-products is 2.5 points below the published result, and the ogbn-products CoE results use a single seed because one run takes about three hours; the gaps we report there (0.5--1.6 points) are large compared with the seed variation of the reference models.
The full-batch GNNs were not used as experts on ogbn-products, so the expert pool there is small.
The cooperation matrix is shared by all nodes; since the value of label propagation depends on how much label information reaches a node, a per-node weighting is a natural next step.

\section{Conclusion}

We evaluated seven reference models on Cora, Citeseer and ogbn-products under one protocol, and adapted Cooperation of Experts to node classification by fitting the cooperation matrix on out-of-fold predictions of experts that use different information.
On ogbn-products, CoE improves over its best expert by 1.6 points and over the best reference model with C\&S by 0.5 points, mainly by learning how to combine label propagation with a feature-based model.
On the citation graphs it only matches the best experts, also with many labels, because the GNN experts make similar errors.
With one round of self-training on the CoE pseudo-labels, ogbn-products accuracy reaches 80.8\%, and SIGN trained on these pseudo-labels is 1.1 points better than on its own.
The large-margin term did not improve over plain stacking, and on the small graphs the gains of ensemble self-training are those of self-training itself.

\bibliographystyle{ACM-Reference-Format}
\bibliography{refs}

\end{document}
